\bookchapter{第零章：嗨，開始之前}

喺正式探索前，我哋先傾啲基本規則，確保後面講嘅你都識。唔好急，磨刀不誤砍柴工嘛。

\booksection{一、呢本書可以幫你做乜嘢？（同唔可以做嘅）}

佢可以幫你係統性咁搞懂手機點樣用，明白功能背後嘅道理，而且無論你係安卓黨定蘋果派，大多數思路係通用嘅。佢會帶住你動手練習，等你由「識用」升級到「懂用」。

佢唔係某一款特定手機嘅說明書，亦唔係追新品嘅科技新聞，更唔係教你沉迷手機嘅指南。講真，如果你期待嘅係「點樣用手機刷影片更爽」，咁可能要失望啦——我哋更關心嘅係點樣等手機真正為你服務，而唔係反過嚟。

\booksection{二、我哋嘅「暗號」說明}

為咗講得更清楚，書裏會用一啲小標記，呢啲標記就似我哋之間嘅默契暗號：

\begin{itemize}
  \item \textbf{[按鈕名]}：代表螢幕上具體嘅按鈕或App名，比如[設定]、[確定]。
  \item \textbf{>}：代表操作順序，比如「撳 [設定] > [聲音]」，就係先撳設定，再撳聲音。
  \item \textbf{加粗}：非常重要！可能係關鍵概念或安全警告，請務必留意。
  \item \emph{斜體}：一啲補充嘅小想法或比喻，幫你更好地理解。
\end{itemize}

\booksection{三、安卓同iOS，我哋都講講}

你估點樣，有啲嘢叫法唔同，但做嘅事一樣。比如：應用（安卓）同 App（iOS），設定（都叫設定），檔案管理（安卓）同 檔案（iOS）。知道係一回事就得，唔好俾名詞繞暈咗。

另外，而家安卓陣營裏仲有華為嘅鴻蒙系統，佢有自己嘅特色，後面我哋會專門講到。總之，唔理你用乜嘢手機，呢本書都幫到你。

\booksection{四、手機容量：點解買咗128G，可用卻唔到110G？}

容量單位大家知道：1TB = 1024GB，1GB = 1024MB，1MB = 1024KB……但呢度有個坑：硬碟廠商通常按1000進制算，而系統按1024進制算，所以你買到嘅「128G」硬碟，實際容量會「縮水」。

但你手機容量縮水嘅原因主要有三個：

\begin{enumerate}
  \item \textbf{系統自己霸咗地方}：作業系統就似房子嘅承重牆，必須存在。而家嘅系統動輒佔用10-20GB，呢個好正常。
  \item \textbf{廠商預裝嘅應用}：新手機裏自帶嘅嗰啲App都佔空間，有啲可以卸載，有啲唔可以。
  \item \textbf{分區同格式化損耗}：技術原因，就似書包嘅襯布都佔體積。
\end{enumerate}

\begin{tcbinfo}{買手機嘅小建議}
而家128GB嘅手機基本唔夠用啦，256GB起步比較穩陣。如果你鍾意影相、玩遊戲，建議直接上512GB甚至更大。記住，存儲空間呢啲嘢，寧可富餘，唔可捉襟見肘。

\end{tcbinfo}

\begin{tcbpractice}{嚟動手：睇下你嘅手機「身份證」}
打開 [設定]，揾到關於手機（或類似選項），記低：手機型號、系統版本、總存儲空間同剩餘空間。呢個係瞭解佢嘅第一步。順便睇下電池健康度，心裏有個數。

\end{tcbpractice}

\booksection{五、你嘅手機係一座「數碼資產庫」}

喺開始具體知識之前，我哋先建立一個核心認知：\textbf{手機早已唔只係通信工具，佢係我哋最重要嘅數碼資產庫}。聯繫人、聊天記錄、家庭相、工作文檔、銀行App、付款碼、健康數據、電子證件……呢啲數據加埋一齊，價值往往超過手機本身。

\begin{tcbinfo}{點解呢個好重要？}
如果手機唔見咗或壞咗，你損失嘅唔只係幾千蚊硬件，\textbf{更可能係多年積累嘅不可再生數據}——小朋友嘅成長相、過世親人嘅語音、冇備份嘅重要合同。科學咁整理同防護習慣，可以從三個維度保護你：

\begin{itemize}
  \item \textbf{提升運行效率}：規範嘅檔案分類等系統查找更快，手機長期唔卡頓。
  \item \textbf{延長使用壽命}：合理嘅存儲管理減少晶片讀寫負擔，電池同存儲都更耐用。
  \item \textbf{保障數據安全}：分層管理 + 定期備份 + 權限收斂，等「數據災難」遠離你。
\end{itemize}

\end{tcbinfo}

\booksubsection{我哋嘅核心理念：「預防為主，防治結合」}

就似養生唔可以等到得病先睇醫生，\textbf{手機管理都要由「日常習慣」抓起}。本書反覆強調嘅唔係某個高級技巧，而係呢啲樸素嘅原則：

\begin{itemize}
  \item \textbf{預防為主}：裝 App 淨係行官方商店、發相之前先脱敏、密碼唔重複——呢啲「小動作」成本極低，但可以擋住絕大多數風險。
  \item \textbf{定期體檢}：每月一次整理下載檔案夾、每季度一次權限審計、每半年一次全量備份——固定節奏比臨時救火有效一百倍。
  \item \textbf{留有退路}：開啟雲同步、記好帳號密碼、寫一份「數碼遺產清單」——確保意外發生時你同屋企人都可以從容應對。
\end{itemize}
